\documentclass[11pt,letterpaper]{article}
\usepackage[T1]{fontenc}
\usepackage[utf8]{inputenc}
\usepackage{lmodern}
\usepackage[margin=1in,headheight=15pt]{geometry}
\usepackage{microtype,amsmath,amssymb,amsthm,mathtools}
\usepackage{booktabs,longtable,array,enumitem,multicol}
\usepackage[dvipsnames]{xcolor}
\usepackage{fancyhdr,hyperref,listings,xurl}
\definecolor{inkblue}{HTML}{17365D}
\definecolor{muted}{HTML}{536471}
\hypersetup{colorlinks=true,linkcolor=inkblue,citecolor=inkblue,urlcolor=inkblue,
 pdftitle={Notations for Omnific Integers: Exact Holonomic Towers and the Complexity of Equality},
 pdfauthor={Research report prepared for Vladimir Reshetnikov}}
\urlstyle{same}
\pagestyle{fancy}\fancyhf{}
\fancyhead[L]{\small\textsc{Notations for omnific integers}}
\fancyhead[R]{\small\textsc{Theory and decision procedures}}
\fancyfoot[C]{\thepage}
\renewcommand{\headrulewidth}{0.3pt}
\setlength{\parskip}{4pt}\setlength{\parindent}{1.15em}
\setlength{\emergencystretch}{2em}
\setcounter{tocdepth}{1}
\hypersetup{bookmarksdepth=2}
\numberwithin{equation}{section}
\newtheorem{theorem}{Theorem}[section]
\newtheorem{lemma}[theorem]{Lemma}
\newtheorem{proposition}[theorem]{Proposition}
\newtheorem{corollary}[theorem]{Corollary}
\theoremstyle{definition}
\newtheorem{definition}[theorem]{Definition}
\newtheorem{example}[theorem]{Example}
\theoremstyle{remark}
\newtheorem{remark}[theorem]{Remark}
\newcommand{\No}{\mathbf{No}}
\newcommand{\Oz}{\mathbf{Oz}}
\newcommand{\On}{\mathbf{On}}
\newcommand{\N}{\mathbb N}
\newcommand{\Z}{\mathbb Z}
\newcommand{\Q}{\mathbb Q}
\newcommand{\R}{\mathbb R}
\newcommand{\DF}{\operatorname{DF}}
\newcommand{\Frac}{\operatorname{Frac}}
\newcommand{\supp}{\operatorname{supp}}
\newcommand{\ct}{\operatorname{ct}}
\newcommand{\lc}{\operatorname{lc}}
\newcommand{\lt}{\operatorname{lt}}
\newcommand{\ord}{\operatorname{ord}}
\newcommand{\Tr}{\operatorname{Tr}}
\newcommand{\sgn}{\operatorname{sgn}}
\newcommand{\WF}{\mathrm{WF}}
\newcommand{\KB}{\mathrm{KB}}
\newcommand{\Pos}{\mathrm{Pos}}
\newcommand{\Eq}{\mathrm{Eq}}
\newcommand{\val}{\operatorname{v}}
\newcommand{\ceil}[1]{\left\lceil #1\right\rceil}
\newcommand{\floor}[1]{\left\lfloor #1\right\rfloor}
\newcolumntype{P}[1]{>{\raggedright\arraybackslash}p{#1}}
\lstset{basicstyle=\small\ttfamily,columns=fullflexible,breaklines=true,
 frame=single,rulecolor=\color{muted},showstringspaces=false}

\begin{document}
\hypersetup{pageanchor=false}
\begin{titlepage}
\vspace*{0.55in}
\noindent{\color{inkblue}\Large\textsc{Foundations, ordinal notation, and effective algebra}}
\vspace{0.32in}

\noindent{\Huge\bfseries Notations for\\[5pt] Omnific Integers}
\vspace{0.25in}

\noindent{\LARGE Exact holonomic towers and\\[4pt] the complexity of equality}
\vspace{0.36in}

\noindent{\large A mathematical article with complete proofs,\\
explicit representation contracts, and a finite verification prototype}
\vspace{0.38in}

\noindent Prepared for Vladimir Reshetnikov\\
23 September 2026

\vfill
\noindent\begin{minipage}{\textwidth}
\textbf{Main results.} A decidable notation system obtained by iterating
certified differential equations has computable equality, order, truncation,
omnific-integer recognition, and omnific floor. With hereditary Cantor-normal-form
scale labels, its ordinal trace is exactly $\varepsilon_0$. For a contrasting
primitive-recursive stream representation, equality is $\Pi^0_1$-complete,
while positivity in $d$ dominance layers is complete at the $d$th finite
mind-change level---already on values with at most $d$ monomials.

\medskip
\textbf{Status.} The proofs in this article are mathematical arguments, not
Lean formalizations. Standard ingredients and repository precedents are
identified explicitly. Several combined results are proposed as original
formulations; the literature search is not a proof of priority. The supplied
program checks a restricted exact-arithmetic implementation, not the infinite
or computability-theoretic theorems.
\end{minipage}
\end{titlepage}
\hypersetup{pageanchor=true}

\begin{abstract}
An omnific integer may have a transfinite Conway normal form, arbitrary real
coefficients on its positive exponents, and an ordinary integer constant
term. Consequently, finite notation is not synonymous with a finite normal
form, and decidable syntax is not synonymous with decidable equality.
We develop three precisely separated types of notation.

First, over an effective ordered real coefficient field, a certified Euler
operator and a finite compatible jet describe a unique D-finite power series.
A leading-polynomial argument computes all exceptional integer indices even
when the coefficients themselves are non-Archimedean. Iterating fraction
fields of these series produces effective finite-rank surreal fields.
Their omnific-integer parts admit a unique triangular decomposition, exact
truncation, a computable floor, and an effective test for being an ordinal.
Varying finitely many scale labels below $\varepsilon_0$ gives one decidable
notation system whose ordinal trace is exactly the ordinals below
$\varepsilon_0$, while also naming nonordinal omnific integers with infinite
support. At rank $r$, support order types are bounded by $\omega^r$; the
strict bound for the integer part is sharp.

Second, primitive-recursive rational coefficient streams on $d$ fixed positive
blocks always denote omnific integers, but their equality relation is
$\Pi^0_1$-complete. Their positivity relation is many-one complete for the
$d$-c.e. sets. The lower bounds require at most one nonzero coefficient per
block. Thus the same very simple values can have decidable or undecidable
notation equality depending on the information supplied by their names.
Third, a fixed computable family of dyadic exponents in $(1,2)$ makes semantic
support validity $\Pi^1_1$-complete, even with all nonzero coefficients equal
to one. These results isolate coefficient, dominance, and well-foundedness
obstructions rather than treating them as one universal undecidability claim.
\end{abstract}

\begingroup\small\setlength{\parskip}{0pt}
\begin{multicols}{2}\tableofcontents\end{multicols}
\endgroup
\newpage

\section{The problem and its three different answers}
\label{sec:intro}

An ordinal notation is a finite object with an ordinal interpretation.
Hereditary Cantor normal forms illustrate one especially favorable situation:
validity and comparison are decidable, and canonical expressions name exactly
the ordinals below $\varepsilon_0$. More general recursive ordinal systems
make well-foundedness part of the semantics, so recognizing valid expressions
can itself become highly undecidable. The overview requested by the user
\cite{ordinalwiki} motivates this distinction; our arguments use the
underlying mathematical facts rather than relying on that overview as a
technical source.

For omnific integers, the analogous question has an additional dimension.
Expressions can contain real coefficients, negative terms, and infinite
normal forms whose exponents remain positive. A system can be very rich in
its values while being algorithmically tame, or extremely restricted in its
values while having an undecidable equality relation. This article makes
both possibilities explicit.

The principal positive construction is denoted $\mathcal R_A$ when its finite
set of scale labels is $A$. It is obtained as the omnific-integer part of an
effective field $\mathcal F_A$. Its names include differential equations with
certified initial data, rational operations, and finitely nested scales.
The principal negative construction uses raw coefficient programs rather
than differential certificates. There, the support geometry is fixed in
advance and harmless, but the name does not certify when a coefficient
stream is identically zero. A third construction releases the support
geometry as well, exposing an independent well-foundedness obstruction.

\begin{center}
\renewcommand{\arraystretch}{1.16}
\begin{tabular}{@{}P{.24\textwidth}P{.23\textwidth}P{.22\textwidth}P{.22\textwidth}@{}}
\toprule
Representation & Semantic validity & Equality & Order / positivity\\
\midrule
Certified holonomic towers & Decidable & Decidable & Decidable\\
$d$ primitive-recursive positive blocks & Every well-formed name is valid &
$\Pi^0_1$-complete & $d$-c.e.-complete positivity\\
Tree-controlled dyadic supports & $\Pi^1_1$-complete & Co-c.e. on valid inputs;
checked equality is $\Pi^1_1$-complete & Not needed for the validity lower bound\\
\bottomrule
\end{tabular}
\end{center}
The last row deliberately distinguishes an equality question under a promise
from a question that must also check the promise.

\subsection{Relation to existing work}

Conway normal forms and the characterization of omnific integers are
classical \cite{conway,gonshor}. Modern omnific-integer work uses generalized
power series in a substantial way; L'Innocente and Mantova develop a
factorisation theory and settle a specific primality conjecture
\cite{innocente-mantova}. Our subject is not factorisation, but finite
representations and their decision problems.

D-finite power-series algebra is classical as well \cite{stanley}. In
particular, closure under addition and multiplication, and finite
initial-value zero tests for specified linear differential equations, are
not new claims here. There are stronger results: Chen, Fang, and van der
Hoeven give a zero-test for D-algebraic transseries in an effective
differential-field framework with specified solutions and transbases
\cite{chen-fang-hoeven}. The construction below is not presented as the
first zero-test beyond rational functions. Its purpose is a transparent
finite certificate format, recursive use of non-Archimedean coefficient
fields, and an exact description of the resulting omnific-integer and ordinal
parts. Our differential operator at a given level treats the lower-level
coefficient field as constant; no compatibility with a global surreal
derivation is assumed.

The inspected Surreal repository already distinguishes structural names,
numerical coefficient names, exact equality, support certificates, and
computability of arithmetic \cite{repo-computable,repo-cas}. Its computer
algebra report contains effective rational ordered-monomial cores and
explicit zero-test and high-rank leading-term obstructions. The present
article does not claim those general principles as new. It adds the
particular certified tower and the exact finite mind-change classification
for omnific positive blocks. Section~\ref{sec:provenance} gives a claim-level
novelty and verification ledger.

\subsection{What ``decidable'' means here}

All algorithms are ordinary terminating algorithms on finite strings, with
no oracle unless explicitly stated. A procedure need not have a practical
complexity bound to be decidable. The coefficient field at the bottom is
part of the input convention, not an arbitrary oracle for real numbers.
The positive construction starts from a fixed effectively presented ordered
subfield $C\subseteq\R$. Its valid codes, equality, order, and field
operations are decidable or computable. Typical choices are $\Q$ and the
real algebraic numbers with exact algebraic representations. Arbitrary
Cauchy names for all computable reals do not satisfy this equality contract.

Every set and sum in the proofs is set-sized. The proper class $\No$ is used
only as an ambient semantic domain. In a universe-relative formalization,
one may choose a universe containing the countable fields and supports under
consideration. No truth predicate for the entire universe is required.

\section{Normal forms and representation contracts}
\label{sec:normal}

\subsection{The semantic domain}

We use Conway's monomial map $x\mapsto\omega^x$, not an unspecified extension
of ordinary exponentiation. Every surreal has a unique normal form
\begin{equation}\label{eq:normal}
 x=\sum_{\gamma\in S} c_\gamma\omega^\gamma,
 \qquad c_\gamma\in\R\setminus\{0\},
\end{equation}
where $S\subseteq\No$ is a set and is reverse well ordered: each nonempty
subset has a greatest element. The sum is a normal-form or Hahn sum.
Its existence is a support assertion, not convergence of a real analytic
series or of the sequence of partial sums in the full surreal order topology.

For $x\ne0$, let $\ell(x)=\max S$ and let $\lc(x)=c_{\ell(x)}$.
The sign of $x$ is the sign of $\lc(x)$. Addition collects equal exponents;
multiplication uses $\omega^\gamma\omega^\eta=\omega^{\gamma+\eta}$
and finite convolution at each exponent. These operations are justified by
the standard support lemmas for generalized power series
\cite{conway,gonshor,innocente-mantova}.

\begin{definition}[Omnific integer]\label{def:oz}
An omnific integer is a surreal of the form
\begin{equation}\label{eq:oz}
 z=m+\sum_{\gamma\in S}c_\gamma\omega^\gamma,
 \qquad m\in\Z,\quad S\subseteq\No_{>0}
\end{equation}
with reverse well-ordered set support. The class of these numbers is $\Oz$.
The positive-exponent coefficients need not be ordinary integers.
\end{definition}

For example, $\sqrt2\,\omega-\tfrac13\omega^{1/2}+7$ is an omnific
integer, whereas $\omega+\tfrac12$ and $\omega-\omega^{-1}$ are not.
The support criterion shows directly that $\Oz$ is a subring of $\No$.
It is discretely ordered: a positive element with a nonempty positive support
is greater than every real number, and otherwise it is a positive ordinary
integer. Thus $1$ is its least positive element.

It is useful to write
\begin{equation}\label{eq:threeparts}
 x=x_{>0}+\ct(x)+x_{<0}.
\end{equation}
The constant term is an ordinary real number. The last part is either zero
or infinitesimal, and the first is either zero or purely infinite. This is a
normal-form decomposition; $x_{>0}$ is not generally the positive part in the
order-theoretic sense.

\subsection{Names, validity, and quotients}

\begin{definition}[Notation system]
A notation system consists of a decidable set $D$ of finite syntactic codes,
a subset $V\subseteq D$ of semantically valid codes, and a denotation map
$\nu:V\to\Oz$. Its equality problem on valid inputs is
$\nu(a)=\nu(b)$ under the promise $a,b\in V$. Its \emph{checked equality}
relation on all syntactic inputs is
\[
 E^{\mathrm{chk}}=\{(a,b)\in D^2:a,b\in V\text{ and }\nu(a)=\nu(b)\}.
\]
A total notation system has $V=D$. A canonicalization is a computable map
$N:V\to V$ with $\nu(N(a))=\nu(a)$ and
$N(a)=N(b)$ exactly when $\nu(a)=\nu(b)$.
\end{definition}

A coefficient-access algorithm answers a different question: it computes a
specified coefficient or a specified finite collection of coefficients.
It does not, by itself, certify that all remaining coefficients vanish.
Likewise, a support generator is not automatically a support-validity
certificate.

\begin{proposition}[Canonicalization and decidable quotients]\label{prop:canonical}
For a total notation system with effectively coded finite strings, equality
is decidable if and only if it has a computable canonicalization.
\end{proposition}
\begin{proof}
A canonicalizer decides equality by comparing its outputs. Conversely,
number all finite strings by natural numbers. For input code $a$, test the
finitely many codes $b\le a$ for syntactic validity and semantic equality to
$a$, and return the least successful one. The set tested is nonempty since
it contains $a$. The result is the least code in the whole equivalence
class: no code larger than $a$ can be the least. Equal inputs therefore give
the same result.
\end{proof}
This proposition supplies a canonical form in the computability sense, not
necessarily a useful algebraic expression or an efficient normalizer.

\begin{proposition}[Cardinality obstruction]
The image of a finite-alphabet notation system is countable. No such system
names all omnific integers, or even all numbers $c\omega$ with $c\in\R$.
\end{proposition}
\begin{proof}
Finite strings form a countable set, whereas $c\mapsto c\omega$ is injective
on $\R$. Moreover $\Oz$ contains every ordinal and is a proper class.
\end{proof}
This elementary obstruction does not prevent large and useful countable
subrings. The question is which semantic closure properties can coexist
with effective equality.

\section{A sparse baseline and the source of additional difficulty}
\label{sec:sparse}

Let $G\subseteq(\No,+,<)$ be an effectively presented ordered abelian group
whose embedding and group operations are fixed. Over the exact coefficient
field $C$, use finite records
\begin{equation}\label{eq:sparse}
 (m;(\gamma_1,c_1),\ldots,(\gamma_s,c_s)),
 \qquad m\in\Z,\quad\gamma_i>0.
\end{equation}
Zero coefficients and repeated exponents may be allowed in raw records.
Sort exponents in decreasing order, collect equal exponents, and remove
zero coefficients.

\begin{proposition}[Sparse normalization]\label{prop:sparse}
The records \eqref{eq:sparse} have decidable validity, equality, and order,
and computable ring operations. After collection and sorting, their list of
exponent--coefficient values and their integer constant term are unique.
\end{proposition}
\begin{proof}
Only finitely many comparisons and coefficient operations are involved.
The normal-form uniqueness theorem proves both equality and uniqueness.
The sign of the first nonzero term decides order; if no positive term remains,
the ordinary integer constant decides it. Products have exponent set in the
finite sumset of the input supports, and positive exponents remain positive.
\end{proof}

If coefficient and exponent codes are not themselves canonical, this is
uniqueness of values rather than literal strings. Proposition~\ref{prop:canonical}
can normalize those codes too. The sparse core is familiar generalized
polynomial algebra and already occurs in the repository's computer algebra
work \cite{repo-cas}.

Adding infinite series is not merely adding longer lists. For instance,
\begin{equation}\label{eq:guardbasic}
 \sum_{n\ge0}a_n\omega^{\omega-n}
 =\omega^\omega\sum_{n\ge0}a_n(\omega^{-1})^n
\end{equation}
always has positive exponents and valid reverse well-ordered support,
whatever the real coefficients $a_n$ are. But an algorithm for each $a_n$
need not decide whether every $a_n$ is zero. A finite differential certificate
can supply the missing information. The next sections make that assertion
precise, including the singular initial indices that cannot simply be ignored.

\section{Computing integer singularities over non-Archimedean fields}
\label{sec:integerroots}

The recursive construction will need to solve a modest but important problem:
given $P(X)\in K[X]\setminus\{0\}$, compute all ordinary nonnegative integers
$n$ for which $P(n)=0$, even though $K$ is non-Archimedean. Bounding roots by
the absolute values of the coefficients in $K$ does not solve this: that
bound may be an infinite surreal with no ordinary integer above it.

\begin{definition}[Effective leading-data field]
An effective leading-data field over $C$ is an effectively presented ordered
subfield $K\subseteq\No$ with support in an effective ordered exponent group
$G$, together with algorithms that, on $a\ne0$, compute
$\ell(a)\in G$ and $\lc(a)\in C$.
\end{definition}
At rank zero, $K=C$, every nonzero element has exponent zero and is its own
leading coefficient.

\begin{lemma}[Archimedean integer search]\label{lem:archsearch}
Given $c\in C$, an ordinary integer greater than $|c|$ can be found
effectively, and $\floor{c}\in\Z$ can be computed.
\end{lemma}
\begin{proof}
Test $2^j>|c|$ for $j=0,1,\ldots$. Archimedeanness guarantees termination.
Within the resulting finite integer interval, exact comparisons determine
the greatest integer at most $c$. Binary search is optional.
\end{proof}

\begin{theorem}[Leading-polynomial integer-root algorithm]\label{thm:integerroots}
Let $K$ be an effective leading-data field over $C$. There is an algorithm
that, given $0\ne P\in K[X]$, returns the complete finite set
\[
 I(P)=\{n\in\N:P(n)=0\}.
\]
In particular, it computes
$B(P)=\max(I(P)\cup\{-1\})$.
\end{theorem}
\begin{proof}
Write $P(X)=\sum_{i=0}^m a_iX^i$, omitting zero coefficients when taking
leading exponents. Put
\[
 g=\max_{a_i\ne0}\ell(a_i),\qquad
 p(X)=\sum_{\ell(a_i)=g}\lc(a_i)X^i\in C[X].
\]
The polynomial $p$ is nonzero: its nonzero contributions have different
powers of $X$, so they cannot cancel as polynomial coefficients.
For each ordinary nonnegative integer $n$, the coefficient of $\omega^g$
in $P(n)$ is exactly $p(n)$. Hence
\begin{equation}\label{eq:rootcontain}
 P(n)=0\quad\Longrightarrow\quad p(n)=0.
\end{equation}
This includes $n=0$; the constant coefficient is evaluated in the usual way.

If $p$ is a nonzero constant, return the empty set. Otherwise, write
$p(X)=b_dX^d+\cdots+b_0$ with $b_d\ne0$. Choose an integer
\[
 M>1+\max_{i<d}|b_i/b_d|
\]
by Lemma~\ref{lem:archsearch}. The elementary Cauchy bound implies that
$p$ has no real root of absolute value at least $M$. To verify the bound
directly, for $|x|>1+\max_{i<d}|b_i/b_d|$, the sum of the lower-degree
terms has absolute value smaller than $|b_dx^d|$, by the geometric-series
estimate. Thus \eqref{eq:rootcontain} reduces the problem to the finite list
$n=0,\ldots,M$. Evaluate $P(n)$ in $K$ and retain exactly its zeros.
\end{proof}

The theorem does not compute all roots in $K$, nor require $K$ to be real
closed. It uses ordinary real leading coefficients to bound possible
\emph{ordinary integer} roots, followed by exact testing in $K$. The leading
polynomial is only a filter: cancellation at its root can leave a nonzero
lower-scale term, which is why the final tests are necessary.

\begin{example}
In a field containing $t=\omega^{-1}$, let
$P(X)=(X-3)\omega+(X-2)$. Its leading polynomial is $X-3$.
The only candidate nonnegative integer root is $3$, but $P(3)=1$, so
$I(P)=\varnothing$. In contrast, $(X-3)(\omega+1)$ has the same leading
polynomial and does have the integer root $3$.
\end{example}

\section{Certified differential-series names}
\label{sec:certificates}

Fix an effective leading-data field $K$ and an independent formal variable
$t$. Differentiation with respect to $t$ kills every element of $K$.
Set $\theta=t\,d/dt$.

\subsection{Euler-normalized operators}

An Euler-normalized operator is
\begin{equation}\label{eq:euleroperator}
 L=\sum_{j=0}^{J}t^jP_j(\theta),\qquad
 P_j(X)\in K[X],\quad P_0\ne0.
\end{equation}
All nonzero linear differential operators over $K(t)$ can be put into this
form without changing their power-series solutions. First clear rational
function denominators. Then use
\begin{equation}\label{eq:falling}
 t^a\left(\frac d{dt}\right)^b
 =t^{a-b}\theta(\theta-1)\cdots(\theta-b+1).
\end{equation}
Collect equal powers of $t$, and multiply by a suitable power of $t$ to make
the least occurring power zero. Multiplication of the equation by a nonzero
rational function does not change its solution set inside $K((t))$.
The uniqueness of the differential-operator normal form, or the triangular
change of basis between powers and falling factorials, ensures that a
nonzero operator stays nonzero.

For $f=\sum_{n\ge0}a_nt^n$, the equation $Lf=0$ is equivalent to
\begin{equation}\label{eq:recurrence}
 P_0(n)a_n+\sum_{j=1}^{J}P_j(n-j)a_{n-j}=0
 \quad(n\ge0),\qquad a_m=0\ (m<0).
\end{equation}
It is important that the argument of $P_j$ is $n-j$, not $n$.

\begin{definition}[Certified holonomic name]\label{def:certificate}
A certified holonomic name over $K$ is an operator \eqref{eq:euleroperator}
together with the finite jet
\[
 (a_0,\ldots,a_B),\qquad B=B(P_0),
\]
where the jet is empty if $B=-1$, and every equation
\eqref{eq:recurrence} with $0\le n\le B$ is satisfied. A syntactically
well-formed but incompatible jet is invalid and is rejected.
\end{definition}

\begin{theorem}[Finite certification, zero testing, and leading data]
\label{thm:certificate}
A certified holonomic name has a unique denotation in $K[[t]]$. Validity,
coefficient extraction, and the zero test are effective. Explicitly,
\begin{equation}\label{eq:jetzero}
 f=0\quad\Longleftrightarrow\quad a_0=\cdots=a_B=0.
\end{equation}
For a nonzero denotation its least nonzero index is the least nonzero
index in the supplied jet, and is at most $B$.
Conversely, every D-finite power series over $K$ has such a name.
\end{theorem}
\begin{proof}
Compute $B$ by Theorem~\ref{thm:integerroots}. There are finitely many
compatibility equations to check. For every $n>B$, $P_0(n)\ne0$, so define
\begin{equation}\label{eq:extendjet}
 a_n=-P_0(n)^{-1}\sum_{j=1}^{J}P_j(n-j)a_{n-j}.
\end{equation}
This is a uniquely determined finite computation from earlier coefficients.
It gives a formal series satisfying every coefficient equation, hence $Lf=0$.
Any solution with that jet must satisfy the same recursion, proving
uniqueness. If the entire jet is zero, induction in \eqref{eq:extendjet}
gives the zero series. The converse of \eqref{eq:jetzero} is immediate.
If $B=-1$, induction starts at $n=0$ and the only solution is zero.

By definition, a D-finite series satisfies a nonzero linear differential
equation over $K(t)$. Normalize that equation as above and supply the finite
jet of this particular series through the computed bound. It satisfies all
compatibility equations and recovers the series uniquely.
\end{proof}

Here ``every D-finite series'' refers to series whose coefficients belong to
$K$, so its finitely many initial coefficients have $K$-names. No algorithm
is asserted for recognizing that an arbitrary coefficient program is
D-finite or for finding its differential equation.

\begin{example}[Why the differential order is not a sufficient jet bound]
\label{ex:resonance}
The operator $L=\theta(\theta-M)$, for an ordinary integer $M>0$, has the
solutions $1$ and $t^M$. Its differential order is two, but $t^M$ agrees with
zero through index $M-1$. The correct bound from $P_0(X)=X(X-M)$ is $B=M$.
A name for $t^M$ must therefore include its exceptional coefficient at $M$.
There is no valid argument that replaces the computed singular-index bound
by the operator order.
\end{example}

\subsection{Closure algorithms, not just closure assertions}

Let $\DF_K[[t]]$ be the ring of D-finite series over $K$. Its familiar
ring-closure theorem \cite{stanley} has the following effective version in
our representation.

\begin{theorem}[Effective addition and multiplication]\label{thm:dfclosure}
From certified names for $f,g\in\DF_K[[t]]$, one can compute certified
names for $f+g$, $-f$, and $fg$. No oracle for algebraic independence or for
minimal differential equations is required.
\end{theorem}
\begin{proof}
For completeness, an explicit linear-algebra construction follows.
A nonzero series satisfying an operator of differential order $r\ge1$
has a derivative vector
\[
 v_f=(f,f',\ldots,f^{(r-1)})^{\mathsf T},\qquad v_f'=A_fv_f,
\]
where $A_f$ is the companion matrix over $K(t)$. The zero case is already
decidable and can be treated separately. For a sum use the direct-sum
system $A_f\oplus A_g$ and an output row selecting $f+g$.
For a product use
\[
 A=A_f\otimes I+I\otimes A_g
\]
on $v_f\otimes v_g$, with the output row selecting $fg$.
These systems have dimensions $r+s$ and $rs$, respectively.

If a scalar output is $h=c_0v$ in a system $v'=Av$, define rows
$c_{i+1}=c_i'+c_iA$. Then $h^{(i)}=c_iv$. In a system of dimension $q$,
the rows $c_0,\ldots,c_q$ are linearly dependent over $K(t)$.
Gaussian elimination computes a nontrivial relation
$\sum_i b_i(t)c_i=0$, yielding the nonzero operator
$\sum_i b_i(t)(d/dt)^i$ annihilating $h$.
All rational-function operations and zero tests reduce to the corresponding
operations in $K$.

Normalize the resulting operator to \eqref{eq:euleroperator}, compute its
bound $B$, and compute the required jet by finite coefficient addition or
Cauchy convolution. The genuine sum or product supplies a solution, so the
jet satisfies the finite compatibility equations. The uniqueness theorem
then verifies the returned name. Negation can retain the original operator
and negate every jet coefficient.
\end{proof}

Repeated operator construction may be expensive and can produce nonminimal
orders. That affects performance, not termination. The proof also explains
why two different differential equations for the same series can be
compared: compute an annihilator and the required finite jet of their
difference, rather than comparing their input operators literally.

\subsection{Passing to a field safely}

D-finite series do not in general form a field. We therefore explicitly take
\begin{equation}\label{eq:H}
 H(K,t)=\Frac(\DF_K[[t]])\subseteq K((t)).
\end{equation}
A name is a pair $(p,q)$ of certified power-series names with $q\ne0$.
The denominator test is decidable by Theorem~\ref{thm:certificate}.

\begin{proposition}[Effective fraction field]\label{prop:fractions}
The field $H(K,t)$ has computable arithmetic and equality. Its Laurent
valuation, leading coefficient in $K$, and individual Laurent coefficients
are computable.
\end{proposition}
\begin{proof}
Use the usual fraction formulas and Theorem~\ref{thm:dfclosure}.
The equality test is
\[
 \frac pq=\frac rs\quad\Longleftrightarrow\quad ps-rq=0.
\]
For a nonzero fraction, write
$p=t^a(p_a+p_{a+1}t+\cdots)$ and
$q=t^b(q_b+q_{b+1}t+\cdots)$ with $p_aq_b\ne0$.
The integers $a,b$ are read from the finite certified jets. Its Laurent
valuation is $a-b$, and its leading coefficient is $p_a/q_b$.
To compute its coefficient at any specified Laurent index, use triangular
series division after removing $t^b$ from the denominator. Every requested
coefficient requires only finitely many earlier coefficients and division
by the fixed nonzero element $q_b$.
\end{proof}
A reciprocal in $H(K,t)$ is represented as a fraction; it is not silently
asserted to be D-finite. Rational functions are included because polynomials
are D-finite. Every D-finite Laurent series also belongs to this field after
multiplication by a sufficiently large power of $t$.

\section{Finite-rank holonomic towers inside the surreal field}
\label{sec:towers}

\subsection{Scales and their order}

Choose a finite increasing list of ordinal labels
\[
 \alpha_1<\cdots<\alpha_r.
\]
Set
\begin{equation}\label{eq:scales}
 \delta_j=\omega^{\alpha_j},\qquad
 t_j=\omega^{-\delta_j},\qquad U_j=t_j^{-1}=\omega^{\delta_j},
 \qquad G_j=\Z\delta_1+\cdots+\Z\delta_j.
\end{equation}
The coefficients in $G_j$ are ordinary integers, and addition in $G_j$ is
surreal addition. The $\delta_j$ are linearly independent over $\Q$ by
normal-form uniqueness. Their integer span is ordered lexicographically
from the \emph{largest} index: the sign of a nonzero vector is the sign of
its last nonzero coordinate. Thus every positive multiple of $\delta_j$
dominates every element of $G_{j-1}$ in absolute value.

The ordinal labels are not exponents in a formal analytic exponential.
They specify actual Conway monomials and consequently a definite embedding
of the exponent lattice in $\No$. Different label choices can have the same
finite-rank algorithmic structure but very different surreal magnitudes.

Define recursively
\begin{equation}\label{eq:tower}
 F_0=C,\qquad F_j=H(F_{j-1},t_j)\quad(1\le j\le r).
\end{equation}
At level $j$, the field $F_{j-1}$ is constant for differentiation in $t_j$.
The recursion on levels is essential; an operator is not allowed to use its
own unknown denotation among its coefficients.

\begin{lemma}[Iterated Laurent embedding]\label{lem:embedding}
There is an injective ordered-field embedding
\[
 C((t_1))\cdots((t_r))\longrightarrow\No
\]
with the substitutions \eqref{eq:scales}. For
$f=\sum_{n\ge N}a_nt_r^n\ne0$, where $a_N\ne0$, its leading data are
\begin{equation}\label{eq:leadingtower}
 \ell(f)=\ell(a_N)-N\delta_r,
 \qquad \lc(f)=\lc(a_N).
\end{equation}
\end{lemma}
\begin{proof}
Proceed by induction on $r$. Expand each coefficient $a_n$ at lower rank.
Every exponent in $a_nt_r^n$ has the form $g-n\delta_r$ with $g\in G_{r-1}$.
If $n<m$, every exponent in the $n$th block exceeds every exponent in the
$m$th block, because $(m-n)\delta_r$ dominates $G_{r-1}$. Each individual
block is reverse well ordered by induction, and the blocks are indexed by
a subset of the integers bounded below. Their ordered union is therefore
reverse well ordered. Distinct integer vectors give distinct surreal
exponents, so no two different blocks overlap.

For multiplication, a prescribed top coordinate in a product receives only
finitely many pairs of top indices, since both Laurent series are bounded
below. Within each such pair, the lower-rank convolution is legitimate by
induction. Hence expansion respects addition and multiplication.
Nonzero series have the leading term in \eqref{eq:leadingtower}, proving
injectivity and order preservation. Inverses are respected since a field
homomorphism sends the algebraic inverse to the inverse. The rank-zero
case is the fixed embedding $C\subseteq\R\subseteq\No$.
\end{proof}

\begin{theorem}[Effective holonomic tower]\label{thm:tower}
For every finite list of labels, \eqref{eq:tower} gives a countable
subfield $F_r\subseteq\No$ with a decidable finite-name representation.
It has computable field operations, equality, order, leading exponent,
real leading coefficient, and coefficient extraction at each specified
exponent of $G_r$.
\end{theorem}
\begin{proof}
At rank zero the hypotheses hold by the contract for $C$. Suppose they hold
at rank $j-1$. Theorem~\ref{thm:integerroots} computes the exceptional
indices for every Euler operator over $F_{j-1}$. Therefore finite
certificates over that field have decidable validity and zero tests.
Theorem~\ref{thm:dfclosure} and Proposition~\ref{prop:fractions} give the
field algorithms at rank $j$.

Lemma~\ref{lem:embedding} supplies its surreal denotation and computes its
leading exponent and real leading coefficient by descending one rank.
The sign is the sign of the latter. To extract the coefficient at
$g=m\delta_j+h$, first extract the Laurent coefficient at index $-m$
and then its coefficient at $h\in G_{j-1}$. This is finite recursion.
Each algorithm calls lower-rank algorithms and finitely many same-rank
finite algebra operations; none calls an unconstructed next rank.
The codes are finite trees, so there are countably many.
\end{proof}

There is no unproved claim of analytic convergence in this construction.
For example, a formal factorial series is accepted just as legitimately as
an exponential series. Nor is real closedness asserted: the tower is a
specific effective subfield, not an algebraic closure or the full field of
all surreal numbers with a prescribed support group.

\subsection{A useful modular generalization}

The role of differential equations can be isolated. Suppose a class
$\mathcal A(K,t)\subseteq K[[t]]$ has decidable finite names, computable
ring operations and coefficients, and a computable finite zero-test bound
for each name. Suppose it contains $K[t]$. Then the same fraction-field and
finite-rank construction applies, provided the resulting fields carry the
leading data needed at the next stage. Certified D-finite series give one
explicit, completely specified realization of these assumptions.

This observation permits later replacement by stronger certified series
classes without changing the omnific-integer arguments below. It does
\emph{not} justify inserting arbitrary differential-equation solvers:
a zero test, a chosen-solution convention, and terminating coefficient
algorithms remain actual hypotheses.

\section{Truncation, triangular integer forms, and floor}
\label{sec:integerpart}

\subsection{Finite descriptions of possibly infinite truncations}

For $g\in G_r$, write $\Tr_{\ge g}(f)$ for the sum of normal-form terms of
$f$ with exponent at least $g$, and similarly define $\Tr_{>g}(f)$.
The resulting truncation need not have finite support.

\begin{theorem}[Effective truncation closure]\label{thm:truncation}
The field $F_r$ is closed effectively under $\Tr_{\ge g}$ and $\Tr_{>g}$
for every $g\in G_r$. In particular, $f_{>0}$, $\ct(f)$, and $f_{<0}$
can be computed as finite names.
\end{theorem}
\begin{proof}
The assertion is immediate at rank zero, where the only exponent is zero.
At rank $r$, take $f=\sum_{n\ge N}a_nt_r^n$ and
$g=m\delta_r+h$ with $h\in G_{r-1}$. For $f=0$ return zero.
Block dominance gives the exact identity
\begin{equation}\label{eq:truncation}
 \Tr_{\ge g}(f)
 =\sum_{N\le n<-m}a_nt_r^n+
 t_r^{-m}\Tr_{\ge h}(a_{-m}),
\end{equation}
where an empty sum is zero and coefficients below $N$ are zero.
The same formula with $>$ in the recursive term gives $\Tr_{>g}(f)$.
The first sum is finite. Each coefficient has an effective $F_{r-1}$-name,
and the recursive term is available by induction. Every term is in $F_r$,
so the formula supplies the required finite name. Constant extraction is
the iterated coefficient at zero; subtraction yields the negative part.
\end{proof}

A subtlety worth retaining in an implementation is that an exponent cutoff
and a number-of-terms cutoff are different interfaces. Formula
\eqref{eq:truncation} returns a finite \emph{expression} for the exact
truncation, which can have support of order type $\omega^{r-1}$ or larger.
It does not promise to print an infinite list as a finite list.

\subsection{The exact omnific-integer part}

Put $R_r=F_r\cap\Oz$. Define $R_0=\Z$ inside $C$.

\begin{theorem}[Triangular omnific normal form]\label{thm:triangular}
For each $r\ge1$,
\begin{equation}\label{eq:Rrecursion}
 R_r=R_{r-1}+U_rF_{r-1}[U_r].
\end{equation}
More explicitly, every element of $R_r$ has a unique finite decomposition
by coefficient values
\begin{equation}\label{eq:triangular}
 z=m+\sum_{j=1}^r\sum_{k=1}^{d_j}U_j^k a_{j,k},
 \qquad m\in\Z,\quad a_{j,k}\in F_{j-1},
\end{equation}
after deleting trailing zero coefficients at each level. Recognition of
$R_r$ inside $F_r$ and computation of this decomposition are decidable.
Arithmetic, equality, and order in this notation are computable.
\end{theorem}
\begin{proof}
Write $f=\sum_{n\ge N}a_nt_r^n$. All exponents of a nonzero coefficient
block with $n<0$ are positive, since $-n\delta_r$ dominates every
lower-rank exponent. All exponents of a nonzero block with $n>0$ are
negative. The block $n=0$ has exactly its lower-rank normal form.
Consequently $f$ is omnific if and only if the positive-$t_r$ tail is zero
and $a_0\in R_{r-1}$. There are only finitely many negative Laurent
indices, and their coefficients may be arbitrary elements of $F_{r-1}$.
This proves \eqref{eq:Rrecursion} inductively.

For the algorithm, extract the finite nonpositive Laurent prefix and compare
it with $f$ using the field equality algorithm. If they differ, the
positive-$t_r$ tail is nonzero and $f\notin R_r$. Otherwise apply the
lower-rank recognition algorithm to $a_0$. At rank zero,
Lemma~\ref{lem:archsearch} computes $\floor{c}$ and decides whether it equals
$c$, which is exactly membership in $\Z$.

Different top indices have disjoint support blocks. Equality of two such
expressions therefore forces equality of each top coefficient and of the
constant block. Descending the ranks proves uniqueness of
\eqref{eq:triangular}. The field algorithms decide arithmetic and order;
the resulting omnific elements are normalized again by the same algorithm.
Closure also follows directly from \eqref{eq:Rrecursion}: products of
positive-degree polynomials retain positive top degree, and $R_{r-1}$ is a
subring of the coefficient field.
\end{proof}

The triangular form is a finite outer normal form with exact field
coefficients; it is not a claim that every coefficient has a finite Conway
normal form. Applying Proposition~\ref{prop:canonical} to coefficient names
would additionally give a literal canonical string, but is not necessary
for an effective equality test.

\begin{corollary}[Normal-form membership test]\label{cor:oztest}
For $x\in F_r$, let $p=x_{>0}$, $c=\ct(x)$, and $\epsilon=x_{<0}$.
Then
\[
 x\in R_r\quad\Longleftrightarrow\quad \epsilon=0\ \text{and}\ c\in\Z.
\]
Both tests are effective. The element $p$ always belongs to $R_r$.
\end{corollary}

\subsection{The infinitesimal correction in floor}

\begin{theorem}[Computable omnific floor]\label{thm:floor}
For every $x\in F_r$ there is a unique $z\in R_r$ with
$z\le x<z+1$, and it is computable. With the decomposition of
Corollary~\ref{cor:oztest}, it is
\begin{equation}\label{eq:floor}
 \floor{x}_{\Oz}
 =p+\floor{c}_{\Z}-
 \begin{cases}
 1,&c\in\Z\text{ and }\epsilon<0,\\
 0,&\text{otherwise}.
 \end{cases}
\end{equation}
\end{theorem}
\begin{proof}
The element on the right is omnific by Corollary~\ref{cor:oztest}.
If $c$ is not an integer, its distances from its two adjacent integers are
strictly positive real numbers; an infinitesimal perturbation cannot cross
either boundary. Thus $p+\floor{c}$ is the desired integer part.
If $c$ is an integer and $\epsilon\ge0$, then $0\le\epsilon<1$.
If $c$ is an integer and $\epsilon<0$, then $0<1+\epsilon<1$.
These are exactly the two remaining cases in \eqref{eq:floor}.
Uniqueness follows because two distinct omnific integers differ in absolute
value by at least one. All terms and tests in the formula are effective.
\end{proof}

\begin{example}
For $x=-\omega+3-\omega^{-1}$ the floor is $-\omega+2$, not $-\omega+3$.
Deleting negative-exponent terms is therefore not an omnific floor algorithm.
For $x=\omega+\tfrac32-\omega^{-1}$ the floor is $\omega+1$; the
nonintegral constant has a genuine real gap to the adjacent integers.
\end{example}

\section{Support order types and the ordinal part}
\label{sec:ordinals}

\subsection{Sharp finite-rank support bounds}

Let $\ord(\supp f)$ denote the ordinal order type of the support when read
in decreasing exponent order, and give the empty support order type zero.
This is a property of the normal form, not the birthday of $f$.

\begin{theorem}[Support spectrum]\label{thm:support}
For $r\ge1$,
\[
 f\in F_r\Longrightarrow\ord(\supp f)\le\omega^r,
 \qquad
 z\in R_r\Longrightarrow\ord(\supp z)<\omega^r.
\]
The field bound is attained by a rational expression. For each
$\eta<\omega^r$, a rational expression in $R_r$ has support order type
exactly $\eta$.
\end{theorem}
\begin{proof}
At rank zero, a nonzero coefficient has a single term. An iterated Laurent
series at rank $r$ has at most countably many top blocks, each of order type
at most $\omega^{r-1}$ by induction. Their ordinal sum has order type at
most $\omega^{r-1}\cdot\omega=\omega^r$. The support of an omnific
integer has only finitely many negative-$t_r$ blocks and then one
$R_{r-1}$ block. For $r=1$ its support is finite. For $r>1$ the bound is
\[
 \omega^{r-1}\cdot d+\beta<\omega^r
 \qquad(d<\omega,\ \beta<\omega^{r-1}).
\]

For sharpness in the field, the rational expression
\[
 G_r=\prod_{j=1}^r(1-t_j)^{-1}
\]
has one nonzero coefficient at every vector in $\N^r$, read in the iterated
lexicographic order. Its support has order type $\omega^r$.
Now write an arbitrary $\eta<\omega^r$ as
\[
 \eta=\omega^{r-1}d_r+\omega^{r-2}d_{r-1}+\cdots+d_1,
 \qquad d_j\in\N.
\]
The element
\begin{equation}\label{eq:realizesupport}
 z_\eta=\sum_{j=1}^r\sum_{k=1}^{d_j}
 U_j^k\prod_{\ell<j}(1-t_\ell)^{-1}
\end{equation}
is omnific by Theorem~\ref{thm:triangular}. Each $(j,k)$ block has
support order type $\omega^{j-1}$. All its coefficients are positive, and
the blocks are disjoint, so none cancels. Higher $j$ blocks precede all
lower $j$ blocks, giving exactly the displayed Cantor normal form of $\eta$.
For $\eta=0$ take the zero expression.
\end{proof}

These bounds will contrast with Section~\ref{sec:raw}: arbitrarily large
finite-rank support order types are compatible with decidable equality,
whereas a raw name with at most one monomial can already have an
undecidable zero test.

\subsection{Recognizing actual ordinals}

There are three operations that must not be conflated: surreal addition,
ordinary ordinal addition, and natural (Hessenberg) ordinal addition.
On ordinal inputs the first agrees with the third, not generally with the
second. For example, in surreal arithmetic $1+\omega=\omega+1$, while
ordinary ordinal addition has $1+_{\mathrm{ord}}\omega=\omega$.
Our notation constructors use field operations; ordinary ordinal operations
can instead be performed on the recognized Cantor normal forms.

\begin{lemma}[Ordinal points of the exponent group]\label{lem:ordinalgroup}
For the scales \eqref{eq:scales},
\[
 G_r\cap\On=\N\delta_1+\cdots+\N\delta_r.
\]
\end{lemma}
\begin{proof}
A nonnegative-integer combination is the Cantor normal form of an ordinal.
Conversely, an element of $G_r$ already has Conway normal form
$\sum_j n_j\omega^{\alpha_j}$ with distinct ordinal exponents.
An ordinal's Conway normal form is its Cantor normal form, whose
coefficients are positive ordinary integers. Uniqueness forces every
nonzero $n_j$ to be positive.
\end{proof}

\begin{theorem}[Ordinal slice and its effective recognition]
\label{thm:ordinalslice}
The ordinals in $F_r$, equivalently in $R_r$, are exactly
\begin{equation}\label{eq:ordinaltrace}
 F_r\cap\On=\N[U_1,\ldots,U_r].
\end{equation}
Given an $F_r$-name, one can decide whether its value is an ordinal and,
in the affirmative case, compute its ordinary Cantor normal form.
For labels $\alpha_j=j-1$, this slice is exactly
$\On_{<\omega^{\omega^r}}$.
\end{theorem}
\begin{proof}
An ordinal has a finite Cantor normal form with ordinal exponents and
positive integer coefficients. Every exponent of an $F_r$ element lies in
$G_r$, so Lemma~\ref{lem:ordinalgroup} implies that such a normal form is a
polynomial in the $U_j$ with natural-number coefficients. Conversely, every
such polynomial is an ordinal by its Cantor normal form.

For recognition, expand in $t_r$. Extract the finite Laurent prefix with
indices at most zero and use equality to test that the positive-index tail
vanishes. If it does not, there is an exponent with negative top coordinate,
which cannot be an ordinal exponent. If the tail vanishes, recursively test
each coefficient of this finite polynomial in $U_r$ for membership in
$\N[U_1,\ldots,U_{r-1}]$. At rank zero the test is membership in $\N$,
which is decidable in $C$. The recursive procedure terminates on a finite
number of coefficients. It yields the finite polynomial and hence the
Cantor normal form: each monomial $\prod_jU_j^{n_j}$ has exponent
$\sum_j n_j\omega^{\alpha_j}$, and these exponents can be sorted exactly.

For consecutive labels, the exponent-group ordinals are exactly those
below $\omega^r$. Finite Cantor normal forms with exponents below
$\omega^r$ give exactly the ordinals below $\omega^{\omega^r}$.
\end{proof}

The recognition algorithm is not an algorithm deciding whether all terms of
an arbitrary recurrence are nonnegative. It uses a zero test for the entire
positive-index tail and then finitely many lower-rank ordinal tests. That
finite structure is why no unproved general positivity procedure is hidden
in the theorem.

\section{One notation system containing all ordinals below epsilon-zero}
\label{sec:epsilon}

\subsection{Canonical labels}

Use the usual hereditary Cantor-normal-form trees: zero is a tree, and a
nonzero tree is a finite list
\[
 (\beta_1,n_1),\ldots,(\beta_s,n_s),
 \qquad \beta_1>\cdots>\beta_s,\quad n_i\in\N_{>0},
\]
whose interpretation is $\omega^{\beta_1}n_1+\cdots+\omega^{\beta_s}n_s$.
The exponents are themselves such trees. Comparison is lexicographic,
recursing into the exponent trees. Finite-tree induction proves termination
and decidable validity. The classical Cantor-normal-form theorem identifies
these canonical trees with exactly $\On_{<\varepsilon_0}$; see also the
formal treatment in \cite{bancerek}. Here $\varepsilon_0$ is the least fixed
point of $\beta\mapsto\omega^\beta$.

For a finite set $A=\{\alpha_1<\cdots<\alpha_r\}$ of these labels, denote
the fields and rings above by $\mathcal F_A$ and $\mathcal R_A$.
A global name is the pair of its finite label set and a valid local name.
Empty label sets give $C$ and $\Z$.

\begin{lemma}[Effective insertion of scales]\label{lem:insert}
If $A\subseteq B$, there are computable denotation-preserving embeddings
of local names for $\mathcal F_A$ into those for $\mathcal F_B$.
These preserve the corresponding omnific-integer parts. The resulting
semantic embeddings commute under successive insertions.
\end{lemma}
\begin{proof}
Insert missing labels in increasing order. At an existing variable, map all
lower-level operator coefficients and jet entries through the already
constructed embedding. A nonzero polynomial stays nonzero, and its
ordinary integer roots are unchanged under an injective field embedding.
Thus its exceptional bound and compatibility equations are preserved.
The induced power series has the mapped coefficients by the uniqueness
recursion. Fraction pairs map in the same way, with nonzero denominators
remaining nonzero.

At an inserted variable, regard every old lower-field element as a constant
series. It has the certificate $\theta f=0$ with the single jet entry
$f(0)$ equal to that element. This supplies the missing tower stage.
Lemma~\ref{lem:embedding} identifies both expressions with the same
surreal number under the fixed scale assignments. Hence different orders
of insertion agree semantically, and omnific membership is preserved.
\end{proof}

\begin{theorem}[Epsilon-labelled decidable omnific notation]
\label{thm:epsilon}
The directed unions
\[
 \mathcal F_\varepsilon=\bigcup_{A\subseteq_{\mathrm{fin}}\varepsilon_0}
 \mathcal F_A,\qquad
 \mathcal R_\varepsilon=\bigcup_{A\subseteq_{\mathrm{fin}}\varepsilon_0}
 \mathcal R_A
\]
have finite-name representations with decidable validity, equality, and
order, computable arithmetic, computable omnific recognition and floor,
and computable truncation at every finitely named exponent in their
exponent groups. Moreover,
\begin{equation}\label{eq:epsilontrace}
 \mathcal F_\varepsilon\cap\On
 =\mathcal R_\varepsilon\cap\On=\On_{<\varepsilon_0}.
\end{equation}
Recognition and conversion of this ordinal slice are effective, and the
whole omnific system admits a computable canonicalization.
\end{theorem}
\begin{proof}
Given two names, compute the sorted union of their finite sets of canonical
labels and embed both names in that common field by
Lemma~\ref{lem:insert}. Apply the local algorithms there. A cutoff exponent
can be handled by the same union operation. Validity is checked by finite
label comparison and the finite recursive certificate checks. Local floor
and recognition preserve denotations under insertion, so their answers are
independent of the chosen common field.

If an element in the union is an ordinal, Theorem~\ref{thm:ordinalslice}
expresses it as a finite natural-coefficient polynomial in finitely many
$\omega^{\omega^\alpha}$ with $\alpha<\varepsilon_0$.
Every exponent in its Cantor normal form is a finite natural sum of
$\omega^\alpha$ with such labels and is below $\varepsilon_0$.
A finite Cantor normal form with exponents below $\varepsilon_0$ is itself
below $\varepsilon_0$.

Conversely, let $\xi<\varepsilon_0$. Its Cantor normal form is finite.
For each exponent $\beta$ in that form, take the finite Cantor normal form
of $\beta$ and collect its exponents as labels into $A$. They are all below
$\varepsilon_0$. Then each $\omega^\beta$ is a monomial in the
$U_\alpha=\omega^{\omega^\alpha}$ for $\alpha\in A$, so $\xi$ belongs
to $\mathcal R_A$. The local recognition procedure converts between these
forms effectively. Finally apply Proposition~\ref{prop:canonical} to the
global total notation system consisting of all locally valid names.
\end{proof}

Here the phrase ``total notation system'' uses the decidable set of locally
valid codes as its syntactic domain. Malformed strings or incompatible
certificates can first be rejected. No noncomputable choice of a valid
representative is hidden in the canonicalizer.

\begin{corollary}[Cofinality and an ambient ordinal bound]\label{cor:ambientbound}
For a nonempty finite label set $A$ with greatest label $\alpha$, the
positive integer powers of $U_\alpha=\omega^{\omega^\alpha}$ are cofinal
in the ordered field $\mathcal F_A$. Its least ordinal upper bound is
\[
 \lambda_A=\omega^{\omega^{\alpha+1}}.
\]
For the global field,
\[
 \mathcal F_\varepsilon\subset(-\varepsilon_0,\varepsilon_0),
\]
and $\varepsilon_0$ is its least ordinal upper bound. These are statements
about ordinal upper bounds, not suprema in $\No$.
\end{corollary}
\begin{proof}
For a nonzero $x\in\mathcal F_A$, its leading exponent belongs to the
integer span of the scale generators. Choose a positive ordinary integer
$N$ with $\ell(x)<N\omega^\alpha$. Normal-form dominance then implies
$|x|<\omega^{N\omega^\alpha}=U_\alpha^N$. These powers are themselves
ordinals in $\mathcal F_A$. Their supremum in the ordinal order is
\[
 \sup_{N<\omega}\omega^{\omega^\alpha\cdot N}
 =\omega^{\omega^\alpha\cdot\omega}
 =\omega^{\omega^{\alpha+1}},
\]
where the middle products are ordinary ordinal products. This proves the
local assertion. For the empty label set, ordinary natural numbers are
cofinal in $C$ and its least ordinal upper bound is $\omega$.
For every $\alpha<\varepsilon_0$ the displayed $\lambda_A$ is below
$\varepsilon_0$, so every global field element has absolute value below
$\varepsilon_0$. Equation~\eqref{eq:epsilontrace} shows that no smaller
ordinal is an upper bound for the entire union.

Indeed $\varepsilon_0/2$ is already a smaller surreal upper bound: every
field element is below some ordinal $\beta<\varepsilon_0$, and the surreal
sum $\beta+\beta$ is still an ordinal below $\varepsilon_0$. Thus
$\beta<\varepsilon_0/2$. This explicitly prevents reading ``least ordinal
upper bound'' as ``surreal supremum.''
\end{proof}

\begin{remark}[Epsilon-zero is a design choice]
This theorem does not say that $\varepsilon_0$ is an absolute ceiling for
effective omnific notation. The chosen labels are canonical hereditary
Cantor trees. A different, independently justified decidable ordinal-label
system can be substituted in the finite-rank construction. One must prove
its validity and comparison properties separately; an unrestricted
well-foundedness promise cannot be smuggled in as decidable syntax.
\end{remark}

\begin{example}[Large magnitudes at small rank]
The two labels $A=\{0,\omega\}$ permit the ordinal
$\omega^{\omega^\omega}+\omega^3+2$ at rank two.
Its magnitude is far beyond the consecutive-label rank-two ordinal cutoff
$\omega^{\omega^2}$, but the number of nested coefficient fields is still
two. Label magnitude, finite rank, support order type, and birthday are
four different measurements.
\end{example}

\section{Worked infinite-support notations and exact identities}
\label{sec:examples}

Take $C=\Q$, consecutive labels $0,1$, and abbreviate
\[
 t=t_1=\omega^{-1},\qquad U=U_2=\omega^\omega.
\]
For every $f\in F_1$, the element $Uf$ is omnific, because all exponents
of $f$ are ordinary integers bounded above and $\omega$ dominates them.
In fact $\Z+UF_1[U]\subseteq R_2$; this is the restricted ring implemented
by the accompanying prototype.

\subsection{A factorial tail with no analytic convergence requirement}

Consider
\begin{equation}\label{eq:factorial}
 A(t)=\sum_{n\ge0}n!\,t^n.
\end{equation}
It has the certificate
\[
 L_A=\theta-t(\theta+1)^2,\qquad a_0=1.
\]
Indeed, $P_0(X)=X$ has the sole nonnegative integer root zero, and
\eqref{eq:recurrence} reads
\[
 n a_n-n^2a_{n-1}=0\quad(n\ge1),
\]
so $a_n=n a_{n-1}=n!$. Thus
\[
 UA(t)=\omega^\omega+\omega^{\omega-1}
       +2\omega^{\omega-2}+6\omega^{\omega-3}+\cdots
\]
is a finitely named omnific integer of infinite support. The fact that
\eqref{eq:factorial} has no positive ordinary analytic radius of convergence
is irrelevant to its formal or surreal denotation. Each coefficient and
each multiplication coefficient is obtained by a finite computation.

\subsection{Exponential notation without a global exponential convention}

Let
\[
 E(t)=\sum_{n\ge0}\frac{t^n}{n!},\qquad
 E_-(t)=\sum_{n\ge0}\frac{(-t)^n}{n!}.
\]
Their certificates are $\theta-t$ and $\theta+t$, each with constant
coefficient one. The formal product $E(t)E_-(t)$ has derivative zero and
constant coefficient one. In characteristic zero, a power series with
zero derivative is constant, so
\begin{equation}\label{eq:expidentity}
 (UE(t))(UE_-(t))=U^2.
\end{equation}
Both factors have infinite normal forms and are omnific. The identity is
proved in formal series and transported to $\No$; it does not require a
choice of a global surreal exponential or a theorem identifying every
appearance of the symbol $e$ in some other notation system.

The closure algorithm independently produces a differential certificate
for the difference in \eqref{eq:expidentity} and checks its finite jet.
It need not recognize the identity from the printed expressions.

\subsection{Cancellations and leading terms}

The geometric series
\[
 G(t)=(1-t)^{-1}=\sum_{n\ge0}t^n
\]
has the certificate $\theta-t(\theta+1)$ with initial coefficient one.
The first two coefficients of $G-E$ vanish and the next is $1/2$, hence
\begin{equation}\label{eq:leadingexample}
 \lt\bigl(U(G(t)-E(t))\bigr)=\tfrac12\omega^{\omega-2}.
\end{equation}
The tower computes this leading term after a finite certified zero test;
it does not trust a fixed arbitrary truncation length. Multiplication by
$U$ places the entire remaining infinite tail at positive exponents.

At a higher rank, cancellation of a whole lower-rank infinite series can
occur in a single top coefficient. The recursive field equality test is
precisely what allows the leading-term algorithm to move past such a zero
coefficient. Merely enumerating the globally expanded normal form would
not provide the same guarantee.

\subsection{Ordinal versus nonordinal outputs}

The polynomial expression
\[
 U(t^{-2}+3)+5t^{-2}+7
 =\omega^{\omega+2}+3\omega^\omega+5\omega^2+7
\]
is an ordinal. The recognition procedure outputs the displayed Cantor
normal form. In contrast, $UA(t)$ is not an ordinal: it has infinitely
many terms, and $\omega-n$ for $n>0$ is not an ordinal, despite being a
positive surreal exponent. The coefficient $n!$ being an ordinary
integer does not remedy either issue.

Similarly, $\sqrt2\,U$ is omnific when $C$ contains $\sqrt2$, but is not
an ordinal. The definition of an omnific integer permits arbitrary real
coefficients on positive exponents; confusing those coefficients with
natural-number Cantor coefficients would incorrectly exclude many elements
of the intended domain.

\section{Primitive-recursive stream names: equality is already hard}
\label{sec:raw}

The positive system supplies a certificate that tells when an entire
coefficient sequence vanishes. We now remove that certificate without
introducing any uncertainty about totality or support validity.

\subsection{An everywhere-valid syntax}

Fix a standard decidable grammar of primitive-recursive functions
$\N\to\N$. A rational-valued program is a triple of such programs
$(p,q,b)$, with value $(p(n)-q(n))/(1+b(n))$. Thus its denominator is
positive by syntax, not by an undecidable semantic promise. Integer-valued
programs are the special case $b=0$. Every code in this grammar is total,
and universal evaluation on a code and an argument is computable.
The universal evaluator need not itself be primitive recursive.

For a fixed $d\ge1$, a raw $d$-block name is a tuple of such rational
coefficient programs $a_1,\ldots,a_d$, with denotation
\begin{equation}\label{eq:raw}
 Z(a_1,\ldots,a_d)
 =\sum_{k=1}^d\sum_{n\ge0}a_k(n)\omega^{k\omega-n}
 =\sum_{k=1}^d U^k\sum_{n\ge0}a_k(n)t^n.
\end{equation}
Here $k\omega$ means the surreal integer multiple of $\omega$, equivalently
the ordinal $\omega\cdot k$, not ordinary ordinal multiplication
$k\cdot\omega$. This convention avoids an otherwise serious collapse of
all the blocks.

\begin{proposition}[Total validity and effective ring operations]
\label{prop:rawvalid}
Every raw $d$-block name denotes an omnific integer. Distinct pairs $(k,n)$
have distinct exponents. Allowing finitely many blocks and an ordinary
integer constant gives a notation ring with computable addition,
negation, and multiplication on codes.
\end{proposition}
\begin{proof}
All exponents $k\omega-n$ are positive. Within each block they are strictly
decreasing with $n$, and every exponent in block $k+1$ is greater than every
exponent in block $k$. A finite union of these blocks is reverse well
ordered, even after zero coefficients are omitted. Distinctness follows
from the linear independence of $1$ and $\omega$ in the exponent group.

Addition and negation operate coefficientwise. In a product, the coefficient
in block $k$ at index $n$ is a finite sum of products with block indices
adding to $k$ and ordinary indices adding to $n$. Bounded sums and products
of primitive-recursive rational functions are primitive recursive. The
finite number of blocks can increase, but remains finite. Integer constants
act in the usual way. All product exponents remain positive except for the
ordinary constant term.
\end{proof}

For the exact positivity classification below, we use the zero-constant
syntax \eqref{eq:raw}. Adding a positive ordinary constant changes the
initial-default bookkeeping of the mind-change bound and is therefore
not silently included in that theorem.

\subsection{The exact equality complexity}

Recall that a relation is $\Pi^0_1$ if it can be expressed by a universal
ordinary-natural-number quantifier over a computable predicate; equivalently,
its complement is computably enumerable. Completeness below is under
computable many-one reductions.

\begin{theorem}[Raw equality and zero testing]\label{thm:raw-equality}
For every fixed $d\ge1$, equality of raw $d$-block notations is
$\Pi^0_1$-complete. The zero test has the same complexity. Hardness holds
on a uniformly generated family whose denotations have at most one nonzero
monomial.
\end{theorem}
\begin{proof}
Normal-form uniqueness and the fixed distinct exponent grid give
\[
 Z(a_1,\ldots,a_d)=Z(b_1,\ldots,b_d)
 \quad\Longleftrightarrow\quad
 \forall n\,\forall k\le d\ [a_k(n)=b_k(n)].
\]
The matrix is computable, proving the upper bound.

Let machine $e$ run on empty input. The predicate
\[
 h_e(n)=
 \begin{cases}
 1,&\text{machine $e$ halts for the first time at step $n$},\\
 0,&\text{otherwise}
 \end{cases}
\]
is primitive recursive, uniformly in $e$, because simulation is bounded
by $n$. Put $a_1=h_e$ and all remaining blocks equal to zero. Then
\begin{equation}\label{eq:spike}
 Z_e=\sum_{n\ge0}h_e(n)\omega^{\omega-n}
 =\begin{cases}
 0,&e\text{ never halts},\\
 \omega^{\omega-N},&e\text{ first halts at step }N.
 \end{cases}
\end{equation}
Thus $Z_e=0$ is exactly nonhalting, a $\Pi^0_1$-complete property.
There is at most one nonzero monomial in every value, proving the final
assertion. Equality with a fixed zero name gives the pair-equality lower
bound as well.
\end{proof}

This is an application of a standard halting reduction, not a new
undecidability principle. Its role is the localization: the syntax is total,
all exponents are positive, the support has at most one term in the hard
family, and every such value is already available in the decidable
holonomic system.

\begin{corollary}[Known leading data do not decide equality]
\label{cor:knownleading}
There is no equality algorithm for a uniformly named family of nonzero
omnific integers having the same leading exponent, leading coefficient,
and ordinary constant term, even if each support has at most two terms.
\end{corollary}
\begin{proof}
Use $Y_e=U^2+Z_e$ and compare with $U^2$. The leading exponent is $2\omega$,
the leading coefficient is one, and the constant term is zero for every
$e$. Equality still holds exactly when $e$ never halts.
\end{proof}

\begin{theorem}[Same values, incompatible effective representations]
\label{thm:notranslator}
Every denotation in the family \eqref{eq:spike} belongs to $R_2$ and has
a finite rational-monomial name. There is nevertheless no algorithm
translating all its raw names into any notation system with decidable
equality, while preserving their denotations and a name for zero.
In particular, there is no such translator into the certified holonomic
system.
\end{theorem}
\begin{proof}
A nonzero value is $Ut^N$, a rational monomial in $R_2$; the zero value is
also there. If a total denotation-preserving translator existed, translate
$Z_e$ and compare the result with the translated zero. This would decide
nonhalting, contradicting Theorem~\ref{thm:raw-equality}.
\end{proof}
The obstruction is not that the values lack exact finite descriptions.
It is the absence of a uniform algorithm extracting those descriptions
from this particular input representation.

\subsection{What finite-support information suffices}

\begin{proposition}[An exact support count is enough]\label{prop:count}
For a raw finite-block name promised to have exactly $s<\omega$ nonzero
coefficients, where $s$ is supplied, an algorithm computes its finite
normal form. Replacing the exact count by the upper bound ``at most one''
is insufficient for an equality algorithm.
\end{proposition}
\begin{proof}
If $s=0$, return zero. Otherwise evaluate coefficients in all blocks in
successive complete prefixes until $s$ nonzero coefficients have been
found. The promise implies that the search terminates and that no other
terms remain. Sort the found exponents and return the list. The failure
for an upper bound is exactly \eqref{eq:spike}.
\end{proof}
This is a promise algorithm: it does not certify the supplied count for
arbitrary raw names. An erroneous count can make it diverge or return an
incorrect truncation, so the count must come from a trusted certificate or
be explicitly treated as a precondition.

\begin{corollary}[No uniform finite observation bound]\label{cor:nojet}
No total computable function of a raw name supplies a finite prefix length
whose being zero in every block guarantees that the entire denotation is
zero. There is also no sound, computably enumerable, complete proof system
for equality of all these raw names.
\end{corollary}
\begin{proof}
A prefix bound would turn the zero test into finitely many coefficient
tests. A sound and complete enumerable equality-proof system would make
the equality relation computably enumerable. Since inequality is already
computably enumerable, dovetailing the two searches would decide equality.
Both conclusions contradict Theorem~\ref{thm:raw-equality}.
\end{proof}

\section{The exact finite mind-change spectrum of positivity}
\label{sec:ershov}

Equality of raw names is controlled by one universal quantifier over their
coefficients. Order has a different obstruction: a later discovery in a
higher block overrides every earlier conclusion based on a lower block.
This gives an exact classification in the finite part of the Ershov
hierarchy, rather than just another proof of undecidability.

\begin{definition}[$d$-c.e. set]
For $d\ge1$, a set $A\subseteq\N$ is $d$-c.e. if there is a total
computable approximation $a(e,s)\in\{0,1\}$ such that $a(e,0)=0$,
$\lim_s a(e,s)=\mathbf1_A(e)$, and for each $e$ the value changes at most
$d$ times as $s$ increases. The $1$-c.e. sets are the computably enumerable
sets, and $2$-c.e. sets are also called d.c.e. sets. The finite levels and
their strictness are standard parts of the mind-change hierarchy
\cite{ershov}.
\end{definition}

\begin{theorem}[Dominance-layer completeness]\label{thm:ershov}
For each $d\ge1$, let $\Pos_d$ be the set of raw $d$-block names
\eqref{eq:raw} denoting a strictly positive omnific integer. Then
$\Pos_d$ is many-one complete for the $d$-c.e. sets.
The hardness reduction can be chosen so that each block has at most one
nonzero coefficient, every such coefficient is $1$ or $-1$, and the total
support has at most $d$ terms. Equality on the unrestricted raw $d$-block
syntax remains $\Pi^0_1$-complete.
\end{theorem}
\begin{proof}[Upper bound]
At stage $s\ge1$, compute the coefficients with $0\le n<s$ in every
block. A block is discovered if this prefix contains a nonzero coefficient.
For any discovered block, its first nonzero coefficient and its sign are
known permanently: all earlier indices were checked in the same prefix.
If no block has been discovered, output zero. Otherwise let $k_s$ be the
largest discovered block index and output one exactly when that block's
first nonzero coefficient is positive. Set the stage-zero output to zero.

The only possible changes arise when a previously undiscovered higher
block replaces the current highest discovered block. The sequence of
record block indices is strictly increasing and has length at most $d$.
Thus the output changes at most $d$ times, including a possible initial
change from zero to one. Every nonzero block is eventually discovered.
If all blocks are zero, the output stays zero. Otherwise the eventual
highest discovered block is the highest genuinely nonzero block, and its
first nonzero coefficient is the leading coefficient of the whole number.
The approximation therefore converges to the indicator of $\Pos_d$.
\end{proof}
\begin{proof}[Lower bound]
Let $A$ be any $d$-c.e. set, with a fixed computable approximation starting
at zero. For each input $e$, run a machine that computes its successive
approximation values and emits a change event whenever the value flips.
The machine has at most $d$ change events. Computing a stage can take an
unbounded but finite time; this causes no problem because the following
predicate uses machine steps, not stages of a possibly non-primitive-recursive
approximation.

Let $\chi_{j,e}(n)$ mean that machine step $n$ is exactly the $j$th change
event. Bounded simulation makes this predicate primitive recursive,
uniformly in $(j,e,n)$. It is one at at most one $n$ for each $j$.
Define raw coefficients
\begin{equation}\label{eq:mindchangecode}
 a_{j,e}(n)=(-1)^{j+1}\chi_{j,e}(n),\qquad 1\le j\le d.
\end{equation}
Their codes can be computed uniformly from $e$ by inserting $e$ as a
constant into the fixed bounded-simulation program.

If there are no changes, the resulting omnific integer is zero and
$e\notin A$. If there are $j\ge1$ changes, its highest occupied block is
$j$, whose sole coefficient has sign $(-1)^{j+1}$. The number is positive
exactly when $j$ is odd. Starting from zero, the final value of the
approximation is one exactly when its number of changes is odd. Therefore
\[
 e\in A\quad\Longleftrightarrow\quad
 Z(a_{1,e},\ldots,a_{d,e})>0.
\]
This is the required many-one reduction, with the asserted finite support
and coefficient restrictions. The equality statement is
Theorem~\ref{thm:raw-equality}.
\end{proof}

\begin{corollary}[Strict order hierarchy with fixed equality complexity]
\label{cor:hierarchy}
One raw positive block has c.e.-complete positivity; two blocks have
d.c.e.-complete positivity; in general $d$ blocks attain the $d$th finite
mind-change level. These positivity levels are strict, although all the
corresponding equality problems have the same $\Pi^0_1$-complete complexity.
\end{corollary}
\begin{proof}
Apply Theorem~\ref{thm:ershov} and the strictness of the finite mind-change
hierarchy \cite{ershov}. If a complete set for level $d+1$ were at level
$d$, closure of $d$-c.e. sets under computable preimages would collapse
those two levels.
\end{proof}

The theorem is insensitive to the times of the individual change events.
Those times determine the finite subtractive offsets $-n$ in the exponents,
but any term in a higher block dominates every term in a lower one.
This is the precise mechanism converting logical mind changes into
surreal dominance changes.

\begin{example}[A sign obstruction under a nonzero promise]
For the halting spike $Z_e$ of \eqref{eq:spike}, put
\[
 W_e=U-U Z_e=U-\sum_{n\ge0}h_e(n)\omega^{2\omega-n}.
\]
Every $W_e$ is nonzero and has at most two monomials. If $e$ does not halt,
$W_e=U>0$. If it halts, the negative term in block two dominates $U$,
so $W_e<0$. Thus a nonzero promise does not make the two-block sign
problem decidable. In contrast, for a single block under a nonzero
promise, scanning until its first nonzero coefficient always terminates
and decides the sign.
\end{example}

\begin{proposition}[Known highest occupied block]\label{prop:highestblock}
For a raw finite-block name supplied with the index of its highest nonzero
block, and the promise that all higher blocks vanish, sign and leading
monomial are computable. Neither this information nor even the full
leading monomial is enough to decide equality in general.
\end{proposition}
\begin{proof}
Scan the specified nonzero block until its first nonzero coefficient occurs.
Its exponent and coefficient give the leading monomial, and its coefficient
gives the sign. The promise makes this correct and terminating. The equality
counterexample is Corollary~\ref{cor:knownleading}.
\end{proof}

An ordinary halting oracle $0'$ decides whether each of the finitely many
blocks is identically zero. It can therefore find the highest nonzero block
and then use Proposition~\ref{prop:highestblock}. This gives a uniform
$0'$ upper bound for order over the union of all finite-block syntaxes.
The one-block equality lower bound shows that this oracle strength is
necessary for a general equality decision. No stronger oracle is concealed
in the finite-level results.

\section{Support validity can be harder than every arithmetic level}
\label{sec:validity}

The raw block construction made support validity automatic. To expose the
separate cost of recognizing a valid support, it suffices to use dyadic
exponents in a fixed small real interval and coefficient one everywhere.
No large ordinal labels, exotic real constants, or varying exponent algebra
are needed.

\subsection{A fixed computable dyadic coding of finite sequences}

Let $\N^{<\N}$ be the set of finite sequences of natural numbers. Its
Kleene--Brouwer order is defined by
\[
 s<_{\KB}t
\]
if $t$ is a proper prefix of $s$, or, at their first differing coordinate,
$s$ has the smaller entry. This is a computable linear order.

\begin{lemma}[Dyadic order reversal]\label{lem:dyadic}
There is a fixed computable injective map
\[
 q:\N^{<\N}\longrightarrow\Q\cap(1,2)
\]
whose values are dyadic and such that
$s<_{\KB}t$ if and only if $q(s)>q(t)$.
\end{lemma}
\begin{proof}
Enumerate finite sequences without repetition in an effective order.
At stage $m$, only finitely many values have been assigned. Compare the
new sequence with each earlier sequence and find its required position in
the reverse Kleene--Brouwer order. Let $l<u$ be the values of its immediate
assigned neighbors in that order, using endpoints $1$ and $2$ when a
neighbor is absent. Assign the midpoint $(l+u)/2$. It is a new dyadic value
in $(1,2)$ and respects all earlier comparisons. To compute $q(s)$, run
this finite construction through the stage where $s$ is enumerated.
Induction on stages proves the claimed equivalence and injectivity.
\end{proof}

Let $T_e\subseteq\N^{<\N}$ range over primitive-recursive trees. A raw
tree-support notation denotes the proposed formal expression
\begin{equation}\label{eq:treeexpression}
 S_e=\sum_{s\in T_e}\omega^{q(s)}.
\end{equation}
The grammar is total: for example, start with a primitive-recursive
predicate and define membership to require that predicate on every prefix,
so that the generated set is always a tree. Its semantic validity means
that the displayed support is reverse well ordered.

\begin{lemma}[Tree well-foundedness and its linearization]
\label{lem:kbtree}
A tree $T\subseteq\N^{<\N}$ has no infinite branch if and only if
$(T,<_{\KB})$ is a well-order.
\end{lemma}
\begin{proof}
An infinite branch gives the strictly descending Kleene--Brouwer chain of
its successively longer prefixes. Conversely, suppose $T$ has no infinite
branch. The proper-extension relation on its nodes is well founded.
Induct on that relation. The subtree rooted at a node $s$ is, in
Kleene--Brouwer order, the ordinal sum in increasing $n$ of the subtrees
rooted at its children $s^\frown n$, followed by the singleton $s$.
By induction the child subtrees are well ordered. A countable ordered sum
of well-orders followed by one element is a well-order. This applies at
the root; the empty tree is harmless.
\end{proof}

The index set of well-founded primitive-recursive trees on
$\N^{<\N}$ is $\Pi^1_1$-complete. This standard effective descriptive-set
theoretic fact is also explicitly recorded in
\cite[Theorem 2.6(g)]{cenzer-marek-remmel}. Here $\Pi^1_1$ means definability
by a universal quantifier over functions $\N\to\N$ followed by an
arithmetic condition. It is not the same class as $\Pi^0_1$.

\begin{theorem}[Localized validity complexity]\label{thm:validity}
Semantic validity of \eqref{eq:treeexpression} is $\Pi^1_1$-complete.
This holds with all nonzero coefficients equal to one and all exponents
in the fixed dyadic subset $q(\N^{<\N})\subset(1,2)$.
Every valid expression denotes an omnific integer.
\end{theorem}
\begin{proof}
By Lemma~\ref{lem:dyadic}, the reverse order on the exponent set of
\eqref{eq:treeexpression} is isomorphic to $(T_e,<_{\KB})$.
Lemma~\ref{lem:kbtree} therefore gives the equivalences
\[
 S_e\text{ has valid support}
 \quad\Longleftrightarrow\quad
 (T_e,<_{\KB})\text{ is well ordered}
 \quad\Longleftrightarrow\quad T_e\text{ is well founded}.
\]
This proves hardness. It also proves the upper bound, or one can write it
directly as the nonexistence of an infinite strictly increasing chain of
active exponents, a $\Pi^1_1$ condition on the computable code.
All exponents are positive and the constant term is zero, so every valid
sum is omnific by Definition~\ref{def:oz}.
\end{proof}

\begin{proposition}[Promised equality versus checked equality]
\label{prop:treeequality}
On valid tree-support names,
\[
 S_e=S_f\quad\Longleftrightarrow\quad T_e=T_f,
\]
so equality under the validity promise has a co-c.e. test for inequality.
The checked equality relation, requiring both validity and equality on all
raw input pairs, is $\Pi^1_1$-complete.
\end{proposition}
\begin{proof}
The coefficient of the distinct exponent $q(s)$ is exactly the indicator
of $s\in T_e$. Normal-form uniqueness yields the equivalence, and a
finite node witnessing different tree membership witnesses inequality.
For checked equality, the conjunction of the two $\Pi^1_1$ validity
conditions and the arithmetic condition $T_e=T_f$ is $\Pi^1_1$.
Its restriction to the diagonal $(e,e)$ is exactly the validity predicate,
which gives hardness.
\end{proof}

\begin{corollary}[No complete enumerable support-certificate system]
No sound computably enumerable family of finite certificates recognizes
exactly all valid tree-support names.
\end{corollary}
\begin{proof}
Existence of an accepted finite certificate would make validity computably
enumerable, contrary to its $\Pi^1_1$-completeness.
\end{proof}

This does not undermine the positive certificate syntax in
Section~\ref{sec:certificates}. That syntax certifies a deliberately
restricted, total family of valid supports. It does not purport to accept
every semantically valid program-generated support.

\section{A practical specification for implementations}
\label{sec:implementation}

\subsection{Data types and trustworthy boundaries}

A full implementation of the positive construction should separate three
data layers. An exact coefficient backend supplies $C$, with decidable
validity, equality, order, and field arithmetic. A scale context stores a
sorted finite list of canonical ordinal labels and represents each lattice
exponent by its integer coordinate vector. A recursive field name is either
a $C$-name or a fraction of certified Euler-series names whose operator and
jet coefficients are field names in the preceding context.

An omnific name may use the triangular form \eqref{eq:triangular} directly.
Its constant is an ordinary integer; each positive outer power is an
ordinary positive integer; and the coefficients are arbitrary valid field
names one level lower. This grammar makes omnific membership automatic.
An importer of an unrestricted field name must instead run the recognition
algorithm of Theorem~\ref{thm:triangular}.

The following pseudocode specifies the critical field-to-integer conversion.
A nonzero field name provides an exact Laurent valuation; zero is handled
before calling that operation.
\begin{lstlisting}
OmnificNormalize(x, level):
    if level == 0:
        m := ordinary_floor(x)
        return Integer(m) if x == m else Reject
    if x == 0:
        return Zero
    N := LaurentValuation(x, top_variable)
    prefix := sum(Coefficient(x,n) * top_variable^n
                  for n from N through 0)
    # The sum is empty when N > 0.
    if x != prefix:
        return Reject
    lower := OmnificNormalize(Coefficient(x,0), level-1)
    if lower is Reject:
        return Reject
    return Triangular(lower, all nonzero coefficients at N <= n < 0)
\end{lstlisting}
Every equality sign in this algorithm invokes the certified field equality
test. Replacing it by a numerical tolerance or by comparison of a fixed
number of coefficients would invalidate the proof.

\subsection{Finite certificates for zero and nonzero answers}

At one level, a zero answer for a certified series consists of its normalized
operator, the verified exceptional-index computation, and the compatible
zero jet. A nonzero answer identifies the first nonzero jet coefficient,
whose own nonzeroness is certified recursively. For fractions, equality
uses a certified numerator cross-difference. At the bottom, exact rational
arithmetic or an exact algebraic-number backend discharges the claims.

The integer-root procedure can expose its finite certificate: a leading
polynomial $p$, an ordinary integer Cauchy bound $M$, and the exact tests of
$P(n)$ for $0\le n\le M$. An independent checker can verify this certificate
without trusting a black-box non-Archimedean root solver. A faster
implementation may factor $p$ first and only test its possible integer
roots, but the exhaustive bounded procedure already proves termination.

These are mathematical certificate specifications. The supplied prototype
uses exact arithmetic but does not emit a proof-assistant object or a
complete independent proof trace for every call.

\subsection{Complexity and resource limits}

The construction has no claimed uniform polynomial-time bound. Differential
orders can grow under tensor products; coefficient expressions can grow
under repeated field operations; a singularity $M$ written in binary can
force a jet of length $M+1$ in the elementary certificate format.
The examples $\theta(\theta-M)$ make this last issue explicit.
A production system could use sparse compatible jets with a proved
completion procedure, share expression subtrees, cache zero tests, and use
specialized annihilator algorithms. Those refinements require their own
correctness arguments.

A resource-limited implementation should return a distinct ``not completed''
status when a permitted computation exceeds a limit. That is not a proof
of inequality, invalidity, or nonexistence. Similarly, failing to import a
raw stream into the holonomic backend says nothing about whether its
\emph{value} has a holonomic or even rational name; Theorem~\ref{thm:notranslator}
shows why these questions cannot be identified.

\subsection{Delivered finite verification code}

The accompanying Python prototype implements certified Euler series over
$\Q$, effective sum and product annihilators by finite differential-system
linear algebra, their fraction field, and the guarded ring
$\Z+UF_1[U]\subset R_2$. It also includes finite checks for the dominance
encoding and a small hereditary ordinal-normal-form module. Its scope is
strictly smaller than the full tower theorem: arbitrary coefficient
backends, general higher-rank operators, the full $R_2$ recognition and floor
algorithms, and all $\varepsilon_0$-labelled field embeddings are mathematical
constructions in this article, not claimed implementations.

The exact test list and actual outcomes are stored in
\path{data/verification.json}. All 41 test groups passed in the supplied run
with Python 3.13.5 and SymPy 1.14.0; the report records the actual runtime
versions as well. The finite combinatorial checks include 521 event schedules,
780 arbitrary-sign block assignments, and all ordered comparisons among
121 dyadic-coded tree nodes. The tests exercise differential recurrence
indices, exceptional jets, infinite-series identities through certificates,
Laurent leading data, and finite instances of the dominance encodings.
Their correctness evidence is executable finite algebra. No finite test
establishes a halting reduction, a completeness theorem, a support theorem
for arbitrary trees, or the correctness of every possible input to the code.

\section{Novelty, repository provenance, and limitations}
\label{sec:provenance}

\subsection{Inspected repository snapshot}

The repository was inspected at the following main-branch commit on
23 September 2026:
\begin{center}
\texttt{a45d72292a46dac13300d395101df188efb84eef}.
\end{center}
The directly relevant sources were the repository overview, the reports
catalogue, the opening representation and provenance sections of
\path{docs/foundations-and-computation/computable-surreals/article.tex},
and the detailed overview of
\path{docs/foundations-and-computation/computer-algebra/README.md}.
The bibliography pins these paths. This is a targeted inspection, not a
claim to have read or verified every report or every Lean declaration in
the repository.

The repository catalogue identifies its reports as AI-assisted research
drafts and separates mathematical exposition, finite verification programs,
and Lean coverage \cite{repo-overview}. That separation is retained here. No theorem of this article is
asserted to have been checked by the repository's axiom audit, and no new
Lean file is supplied. We use classical normal-form semantics directly,
not a claim that every semantic bridge in the repository is already
formalized.

\subsection{Claim-level originality ledger}

\begin{longtable}{@{}P{.30\textwidth}P{.62\textwidth}@{}}
\toprule
Result or ingredient & Attribution and status\\
\midrule
\endfirsthead
\toprule Result or ingredient & Attribution and status\\\midrule
\endhead
Conway normal forms; omnific support criterion; Hahn arithmetic & Classical
surreal and generalized-series theory. No novelty claim.\\
Sparse exact monomial algebra; representation-sensitive obstructions &
Classical in substance and already developed in the inspected repository.
Used as a baseline, not presented as a new result.\\
D-finite ring closure and linear-equation zero tests & Classical.
The article gives explicit proofs to expose the required certificate
information. More general transseries zero tests are acknowledged.\\
Leading-polynomial integer-root procedure and recursive certificate format &
An explicit elementary construction here. Its ingredients are standard;
no separate broad priority claim is made for a root-bound observation.\\
Holonomic fraction towers with triangular $\Oz$-part, exact floor,
truncation, ordinal recognition, and labelled union & Proposed original
combined formulation, with proofs in Sections~\ref{sec:towers}--\ref{sec:epsilon}.
The novelty sought is the simultaneous explicit representation contract and
integer/ordinal analysis, not D-finite closure alone.\\
Raw equality hardness & Standard halting method, localized here to valid
positive-support omnific names and then to at most one monomial. The
same-value nontranslation theorem makes the representation distinction
explicit.\\
$d$-layer positivity is $d$-c.e.-complete, already on at most $d$ monomials &
Proposed original exact classification in the omnific-notation setting.
The Ershov hierarchy and the reduction techniques are established theory;
the dominance-layer realization and matching upper bound are proved here.\\
Dyadic support validity and checked equality & Standard well-founded-tree
complexity transferred by an explicit fixed dyadic order embedding.
A localized formulation, not a new discovery of $\Pi^1_1$-completeness of
well-foundedness.\\
\bottomrule
\end{longtable}

The literature search included exact combinations of omnific integers,
notation equality, holonomic or D-finite representations, and finite
mind-change or Ershov levels, as well as current transseries zero-testing
work. No matching published statement of the two proposed combined
packages was located. This supports only the cautious phrase ``to our
knowledge''; it cannot certify absence from all published literature,
unpublished work, or uninspected repository companions.

\subsection{Limits that are part of the results}

The construction is not a notation system for every omnific integer. It
requires an exact coefficient field, finitely many scales per name, and a
specified class of finite differential certificates. It is not closed under
all conceivable surreal analytic functions or arbitrary cuts, and no
real-closedness theorem is claimed for its field. A general routine deciding
whether a raw coefficient program admits one of these certificates is
neither supplied nor implied. The nontranslation theorem already forbids
such an unrestricted uniform importer for some elementary values.

The strict support bounds concern finite rank, not birthdays. The
$\varepsilon_0$ ordinal trace concerns the chosen label syntax, not the
largest possible ordinal in every decidable notation system. The
$d$-c.e. classification concerns zero-constant raw $d$-block names with
primitive-recursive rational coefficients and ordinary finite block count.
The $\Pi^1_1$ result concerns semantic well-foundedness of a different raw
syntax. Moving any of these conclusions to a different representation
requires a new proof.

\section{Further questions suggested by the theory}
\label{sec:future}

The results leave several concrete extensions rather than a single vague
request for a universal surreal simplifier.

\paragraph{More compact certificates.}
The full jet through the last exceptional index can be very long.
A compressed certificate could list free singular coefficients and prove
that all nonsingular coefficients between them are generated by the
recurrence. Developing such a format with a small independently checkable
proof object could substantially improve the explicit algorithms without
altering their denotations.

\paragraph{Stronger effective coefficient towers.}
One can investigate replacing the D-finite numerator and denominator
class by richer uniquely specified differential-algebraic classes. The
current D-algebraic transseries literature \cite{chen-fang-hoeven} gives
relevant tools, but the integer-part construction needs a precise common
field, leading-data access, truncation closure, and compatible parameter
semantics. Those obligations should be verified separately rather than
inferred from the phrase ``zero-test'' alone.

\paragraph{Beyond finitely many dominance priorities.}
Theorem~\ref{thm:ershov} suggests assigning higher effective difference
hierarchies to appropriately well-founded systems of dominance priorities.
The finite proof relies on a computable uniform bound on the number of
record priorities; an infinite or ordinal-indexed version must specify a
validity certificate and an effective descending control measure. This
article does not assert that the finite classification automatically
extends to arbitrary ordinal priorities.

\paragraph{Intrinsic versus presentation-based invariants.}
A finite value can carry an arbitrarily delayed coefficient discovery in
a raw name. It would be useful to classify representation morphisms that
preserve leading-term certificates or exact support counts, and to determine
which natural subclasses admit effective translations. Theorem~\ref{thm:notranslator}
shows that an extensional inclusion of denotation sets alone is much too
weak to answer this question.

\paragraph{Formal verification.}
A manageable formalization target is the Euler recurrence theorem over an
abstract effective field, followed by the leading-polynomial integer-root
lemma and the triangular integer-part theorem over iterated Laurent fields.
The final bridge to actual surreal monomials can then be kept as a separate
semantic layer. This would parallel the repository's separation of finite
algebra, Hahn support, and surreal interpretation without presuming that a
prototype execution has already supplied the formal proof.

\section{Conclusion}

A theory of omnific notation should state what information a name supplies,
not only what number it denotes. Certified differential equations yield
finite equality tests strong enough to support recursive non-Archimedean
coefficient fields, exact omnific truncation and floor, and a decidable
ordinal slice reaching every ordinal below $\varepsilon_0$. None of this
requires flattening an infinite normal form into a finite term list.

Raw coefficient programs on the very same elementary scales behave
differently. Equality is already co-c.e.-complete on a one-monomial family,
and $d$ dominance layers realize exactly $d$ finite mind changes in the
positivity problem. Releasing support validity exposes the still different
$\Pi^1_1$ well-foundedness barrier, even for dyadic exponents in $(1,2)$.
The positive and negative results are compatible because their
representations carry different certificates. That distinction is the
organizing principle of the theory.

\appendix
\section{Dependency map for the principal proofs}

The positive construction uses the following order of dependencies:
\[
\begin{gathered}
 \text{exact ordered real coefficients}
 \Longrightarrow\text{leading-polynomial integer-root bounds}\\
 \Longrightarrow\text{finite Euler certificates}
 \Longrightarrow\text{effective D-finite ring operations}\\
 \Longrightarrow\text{fraction fields and leading data}
 \Longrightarrow\text{finite-rank recursion}\\
 \Longrightarrow\text{truncation and triangular integer parts}
 \Longrightarrow\text{floor and ordinal recognition}\\
 \Longrightarrow\text{effective insertion of finitely many labels}.
\end{gathered}
\]
The field at the next rank is never used to justify the preceding rank's
zero test. The polynomial root procedure needs only real leading
coefficients, not roots or quantifier elimination in the next field.

The negative results are independent of the differential construction.
They use normal-form uniqueness, fixed positive support blocks, bounded
machine simulation, and---only for the support-validity theorem---the
standard completeness of well-founded primitive-recursive trees.
The same-value nontranslation theorem joins the two tracks by observing
that every value in the spike family has a rational monomial name.

\section{Notation reference}

\begin{longtable}{@{}P{.22\textwidth}P{.70\textwidth}@{}}
\toprule Symbol & Meaning\\\midrule\endfirsthead
\toprule Symbol & Meaning\\\midrule\endhead
$\No,\Oz,\On$ & Surreal numbers, omnific integers, and ordinals.\\
$C$ & Fixed effective ordered subfield of the ordinary real numbers.\\
$\ell(x),\lc(x)$ & Greatest normal-form exponent and its real coefficient,
for $x\ne0$.\\
$\ct(x)$ & Coefficient of the exponent zero.\\
$\theta$ & Euler operator $t\,d/dt$ with the preceding coefficient field
held constant.\\
$B(P_0)$ & Largest nonnegative integer root of $P_0$, or $-1$ if none exists.\\
$H(K,t)$ & Fraction field of the D-finite power series in $K[[t]]$.\\
$\alpha_j,\delta_j$ & Ordinal scale label and the exponent-group generator
$\delta_j=\omega^{\alpha_j}$.\\
$t_j,U_j$ & Actual surreal monomials $\omega^{-\delta_j}$ and
$\omega^{\delta_j}$.\\
$G_r$ & Integer span of $\delta_1,\ldots,\delta_r$, ordered with the
largest index most significant.\\
$F_r,R_r$ & Iterated holonomic fraction field and its omnific-integer part.\\
$\mathcal F_A,\mathcal R_A$ & The same objects with their finite label set
$A$ made explicit.\\
$\mathcal R_\varepsilon$ & Directed union over finite label sets below
$\varepsilon_0$.\\
$\ord(\supp x)$ & Order type of support read in decreasing exponent order;
not a birthday.\\
$Z(a_1,\ldots,a_d)$ & Raw positive-block expression
$\sum_{k,n}a_k(n)\omega^{k\omega-n}$.\\
$d$-c.e. & Admits a computable binary approximation starting at zero with
at most $d$ changes on each input.\\
$<_{\KB}$ & Kleene--Brouwer order on finite sequences of natural numbers.\\
\bottomrule
\end{longtable}

\begin{thebibliography}{99}
\bibitem{conway}
J. H. Conway, \emph{On Numbers and Games}, second edition,
A K Peters, 2001. Originally published in 1976.

\bibitem{gonshor}
H. Gonshor, \emph{An Introduction to the Theory of Surreal Numbers},
London Mathematical Society Lecture Note Series 110,
Cambridge University Press, 1986.

\bibitem{innocente-mantova}
S. L'Innocente and V. Mantova,
\emph{A factorisation theory for generalised power series and omnific integers},
Advances in Mathematics (2024), article 109513.
\href{https://doi.org/10.1016/j.aim.2024.109513}{doi:10.1016/j.aim.2024.109513}.
Preprint \href{https://arxiv.org/abs/1710.07304v5}{arXiv:1710.07304v5}, 2024.

\bibitem{stanley}
R. P. Stanley, \emph{Differentiably finite power series},
European Journal of Combinatorics \textbf{1} (1980), 175--188.
\url{https://math.mit.edu/~rstan/pubs/pubfiles/45.pdf}.

\bibitem{chen-fang-hoeven}
S. Chen, H. Fang, and J. van der Hoeven,
\emph{A zero-test for D-algebraic transseries}, 2026.
\href{https://arxiv.org/abs/2602.01188}{arXiv:2602.01188}.
Author-hosted version dated 6 February 2026:
\url{https://www.texmacs.org/joris/zttr/zttr.html}.

\bibitem{ershov}
Yu. L. Ershov, \emph{A hierarchy of sets. I},
Algebra and Logic \textbf{7} (1968), 25--43.
\href{https://doi.org/10.1007/BF02218750}{doi:10.1007/BF02218750}.

\bibitem{cenzer-marek-remmel}
D. Cenzer, V. W. Marek, and J. B. Remmel,
\emph{Index sets for finite normal predicate logic programs}, 2013.
\href{https://arxiv.org/abs/1303.6555}{arXiv:1303.6555}.
See Section 2, especially Theorem 2.6(g), for the well-founded-tree
index-set fact used here.

\bibitem{bancerek}
G. Bancerek, \emph{Epsilon numbers and Cantor normal form},
Formalized Mathematics (2009).
\href{https://doi.org/10.2478/v10037-009-0032-8}{doi:10.2478/v10037-009-0032-8}.

\bibitem{ordinalwiki}
Wikipedia contributors, \emph{Ordinal notation}, accessed 23 September 2026.
\url{https://en.wikipedia.org/wiki/Ordinal_notation}.
User-specified introductory reference, not a substitute for the proofs
or primary mathematical sources.

\bibitem{repo-computable}
V. Reshetnikov, Surreal project, \emph{Computable Surreal Numbers:
Structural names, Puiseux series, and effective left-finite fields},
reconciled research report, 21 September 2026.
Inspected at commit \texttt{a45d72292a46dac13300d395101df188efb84eef}.
\href{https://github.com/VladimirReshetnikov/Surreal/blob/a45d72292a46dac13300d395101df188efb84eef/docs/foundations-and-computation/computable-surreals/article.tex}{Pinned LaTeX source in the Surreal repository}.

\bibitem{repo-cas}
V. Reshetnikov, Surreal project,
\emph{Surreal and Surcomplex Numbers in Symbolic Computer Algebra},
research-report overview and implementation-scope record.
Inspected at the same commit.
\href{https://github.com/VladimirReshetnikov/Surreal/blob/a45d72292a46dac13300d395101df188efb84eef/docs/foundations-and-computation/computer-algebra/README.md}{Pinned report overview in the Surreal repository}.

\bibitem{repo-overview}
V. Reshetnikov, Surreal project,
\emph{The surreal and surcomplex reports}, catalogue and provenance overview.
\href{https://github.com/VladimirReshetnikov/Surreal/blob/a45d72292a46dac13300d395101df188efb84eef/docs/README.md}{Pinned catalogue at the inspected commit},
accessed 23 September 2026.
\end{thebibliography}
\end{document}
